\section*{Chapitre \ref{ch:mx:request-for-quote-and-dealer-markets-in-theory} --- Demande de cotation et marchés de dealers en théorie}

\begin{solution}{exo:mx:request-for-quote-and-dealer-markets-in-theory:1}
$\mu_{lo}=\lambda_ds/(\lambda_u+\lambda_d+\rho)=0.2\times0.8/102.2=0.0016$.
\end{solution}

\begin{solution}{exo:mx:request-for-quote-and-dealer-markets-in-theory:2}
$C(1)=10$ (sans concurrence, sans fuite)~; $C(2)=10-(\sqrt{29}-2)\times0.564+1=9.09$ points de base.
\end{solution}

\begin{solution}{exo:mx:request-for-quote-and-dealer-markets-in-theory:3}
Lorsqu'elle égale la sélection adverse, 1.5 point de base~: les deux clients sont alors cotés au coût du dealer, 2.0.
\end{solution}

\begin{solution}{exo:mx:request-for-quote-and-dealer-markets-in-theory:4}
Avec $z=0$, le dealer rétrocède tout le surplus à l'investisseur~: le cours acheteur et le cours vendeur sont égaux au prix interdealers. La fourchette vaut $z(\Delta V_h-\Delta V_l)$,
et l'écart entre les valeurs de réserve des deux types se réduit lorsqu'un type faible peut vendre plus tôt~: il supporte le coût de détention pendant moins longtemps.
\end{solution}

\begin{solution}{exo:mx:request-for-quote-and-dealer-markets-in-theory:5}
La composante commune accroît la dispersion des cours acheteurs, mais chaque dealer la décote en conséquence, de sorte qu'elle ajoute un bruit exempt de fuite
sans faire autant baisser le prix gagnant espéré~: le gain par dealer supplémentaire diminue, alors que la fuite par dealer ne diminue pas.
\end{solution}

\begin{solution}{exo:mx:request-for-quote-and-dealer-markets-in-theory:6}
Un dealer périphérique n'atteint la plupart des vendeurs que par le cœur, de sorte que l'obligation transite par davantage de dealers (2.58 étapes en moyenne, contre 1.77) et que chacun
prélève une marge. Des liens vers davantage de dealers du cœur, ou des liens directs entre dealers périphériques, raccourciraient les chaînes.
\end{solution}

\begin{solution}{exo:mx:request-for-quote-and-dealer-markets-in-theory:7}
Trois dealers, 7.77 points de base (contre deux dealers et 9.09 avec $\sigma_c=2$).
\end{solution}

\begin{solution}{exo:mx:request-for-quote-and-dealer-markets-in-theory:8}
Chaque dealer consulté et non retenu apprend qu'un vendeur est sur le marché, ce qui coûte au client une fuite~: au-delà de quelques dealers, la fuite croît plus vite que
l'avantage que procure la concurrence (à un demi-point de base par dealer perdant, quatre est optimal et douze coûtent 2 points de base de plus).
\end{solution}

\begin{solution}{pb:mx:request-for-quote-and-dealer-markets-in-theory:1}
\textbf{1.} Voir la définition.
\textbf{2.} Les types élevé et faible basculent aux taux $\lambda_u$ et $\lambda_d$~; les détenteurs de type faible vendent et les non-détenteurs de type élevé achètent par l'intermédiaire de dealers rencontrés au taux $\rho$~;
les dealers conservent une part $z$ du surplus.
\textbf{3.} $\lambda_d\mu_{ho}=(\lambda_u+\rho)\mu_{lo}$ avec $\mu_{ho}+\mu_{lo}=s$~: $\mu_{lo}=\lambda_ds/(\lambda_u+\lambda_d+\rho)$.
\textbf{4.} Voir la section 1~; cours vendeur $\Delta V_h$, cours acheteur $z\Delta V_l+(1-z)\Delta V_h$.
\textbf{5.} La fourchette baisse~: 210 points de base à 10 rencontres par an, 36.7 à 100, 4.0 à 1\,000.
\textbf{6.} Les dealers soumettent chacun un seul prix sans voir ceux des autres, et le meilleur prix l'emporte à son propre prix.
\textbf{7.} Le gagnant est le dealer qui a l'estimation la plus élevée de la valeur commune~; sa surestimation espérée vaut $\sigma_c\E[\max\text{ de }n]$.
\textbf{8.} $C(n)=m-(\sqrt{\sigma_c^2+\sigma_p^2}-\sigma_c)\E[\max]+\ell(n-1)$~: la marge, la concurrence entre valeurs privées, la fuite.
\textbf{9.} À moins de 0.001 point de base près pour cinq dealers.
\textbf{10.} \emph{Résultat nommé}~: 16 dealers pour une fuite de 0.1 point de base par dealer perdant, 7 pour 0.25, 4 pour 0.5, 2 pour 1, et un seul dealer à partir de 2
points de base, avec des coûts espérés de 5.52, 6.92, 8.02, 9.09 et 10.00.
\textbf{11.} Pour une fuite de 0.5~: 6 dealers sans erreur sur la valeur commune, 4 avec $\sigma_c=2$, 3 avec $\sigma_c=5$.
\textbf{12.} Un client type sollicite peu de dealers~; une exposition plus large est atténuée par la malédiction du vainqueur et par les relations.
\textbf{13.} Seulement si la fuite par dealer perdant est voisine d'un quart de point de base~; il faut la mesurer avant de décider.
\textbf{14.} Quelques pôles parmi de nombreux dealers périphériques~; les obligations affluent vers le cœur~; les dealers centraux facturent des marges aller-retour jusqu'au double, mais offrent
l'immédiateté.
\textbf{15.} Les clients des dealers périphériques empruntent des chaînes plus longues (2.58 étapes contre 1.77) et paient 10.4 points de base contre 9.2.
\textbf{16.} Voir la définition~; avec les chiffres du chapitre, 3.5 à 0.5 point de base pour le client informé à mesure que la valeur de son information passe de 0 à
3, contre 2.0 pour le client non informé.
\textbf{17.} Les prix s'améliorent avec le nombre de dealers jusqu'à une vingtaine, puis se dégradent au-delà.
\textbf{18.} Les dealers pratiquent des fourchettes plus faibles envers leurs contreparties les plus proches, d'autant plus en période de turbulences~; les dealers d'importance systémique facturent plus cher la périphérie et
davantage encore les clients.
\textbf{19.} Comparer l'évolution du prix après les demandes de cotation (RFQ) avec beaucoup ou peu de dealers, ou après les RFQ qui n'ont pas abouti à une transaction, à des obligations comparables sans RFQ.
\textbf{20.} La concurrence apportée par un cours acheteur de plus, face à l'information qu'emporte un perdant de plus.
\end{solution}

\begin{solution}{iq:mx:request-for-quote-and-dealer-markets-in-theory:1}
Peu~: deux à quatre, choisis pour leurs axes et leur historique de cotation, car chaque dealer perdant apprend qu'il existe un vendeur sur une obligation qui s'échange rarement~; davantage lorsque
la fuite est faible (l'obligation est liquide) ou la dispersion des valeurs privées est élevée.
\iqlookfor{Concurrence contre fuite.}
\end{solution}

\begin{solution}{iq:mx:request-for-quote-and-dealer-markets-in-theory:2}
Gagner signifie avoir l'estimation la plus optimiste de la valeur commune~; un dealer doit proposer son estimation diminuée de la surestimation espérée conditionnellement
à la victoire, qui croît avec le nombre de concurrents et la dispersion des estimations.
\iqlookfor{Conditionnement sur la victoire~; la décote croît avec n.}
\end{solution}

\begin{solution}{iq:mx:request-for-quote-and-dealer-markets-in-theory:3}
Le surplus sur lequel négocie un dealer correspond au coût de détention de l'actif supporté par le type faible pendant son attente~; un contact plus facile raccourcit cette attente, réduisant
l'écart entre les valeurs de réserve des types, et donc la fourchette.
\iqlookfor{Valeurs de réserve et attente.}
\end{solution}

\begin{solution}{iq:mx:request-for-quote-and-dealer-markets-in-theory:4}
Parce que ce qu'il apprend du flux vaut davantage dans ses autres transactions que la sélection adverse sur celle-ci~: c'est la \omterm{def:mx:request-for-quote-and-dealer-markets-in-theory:chasing}{chasse à l'information}.
\iqlookfor{Valeur de l'information contre sélection adverse.}
\end{solution}

\begin{solution}{iq:mx:request-for-quote-and-dealer-markets-in-theory:5}
Chaque RFQ, avec les dealers sollicités, leurs cotations et leurs horodatages, le gagnant, ainsi que les cours et transactions du marché sur l'obligation et les obligations liées, avant et
après~; de sorte que les variations de prix après les RFQ puissent être comparées selon le nombre de dealers.
\iqlookfor{Données au niveau de la RFQ~; cours avant et après.}
\end{solution}

\begin{solution}{iq:mx:request-for-quote-and-dealer-markets-in-theory:6}
Les chaînes s'allongent et les marges augmentent lorsque les dealers périphériques contournent le cœur affaibli, et les clients qui ont besoin d'immédiateté paient davantage~; après 2008, les chaînes
se sont allongées de 20~\% (Di Maggio, Kermani et Song).
\iqlookfor{Chaînes, marges, immédiateté.}
\end{solution}
